\clearpage
\fancyhead[L]{\small PIAZZI. STERNVERZEICHNIS.}
\hypertarget{native-595-lower}{}
\begin{center}\small Göttingische gelehrte Anzeigen. Stück 91. Seiten 903..904. 10. Juni 1815.\\[3pt]\rule{.72\linewidth}{.3pt}\end{center}

\emph{Praecipuarum stellarum inerrantium positiones mediae ineunte saeculo XIX ex observationibus habitis in specula Panormitana ab anno} 1792 \emph{ad annum} 1813. \emph{Ex regia typographia militari.} Palermo 1814. 187 Seiten in Folio. (Preis in Florenz bei Molini, Landi und Co. 30 Lire.)

Seit dem Erscheinen des berühmten Sternverzeichnisses von \textsc{Piazzi} aus dem Jahr 1803, der Frucht zehnjähriger ununterbrochener und angestrengter Arbeit, waren die Beobachtungen der Fixsterne an der Sternwarte von Palermo stetig fortgesetzt worden, teils von \textsc{Piazzi} selbst, teils von seinem geschickten Gehilfen \textsc{Nicolao Cacciatore}. Sie erweiterten die Beobachtungen und stellten sie sogar auf eine neue Grundlage, indem sie die Fundamentalrektaszensionen von $\alpha$ im Adler und $\alpha$ im Kleinen Hund durch unmittelbaren Vergleich mit der Sonne bestimmten. Der vorliegende neue Katalog enthält nun den vollständigen Ertrag der gesamten zwanzigjährigen Arbeit. Er gibt die Positionen von 7646 Fixsternen mit der ganzen Genauigkeit an, die die beobachtende Astronomie heute erreichen kann.

In diesem klassischen Werk hat \textsc{Piazzi} auch der Eigenbewegung der Sterne besondere Sorgfalt gewidmet. Er stützte sich dabei vor allem auf die Beobachtungen von \textsc{Mayer} und \textsc{Bradley} und zog sehr häufig auch den Vergleich seiner eigenen Beobach\SourceLeaf{596}tungen untereinander heran. Die Positionen von 1041 Sternen für 1756 sind am Schluss des Werks beigefügt, so wie sie ein anderer Gehilfe von \textsc{Piazzi}, \textsc{Joseph Pilati}, zu diesem Zweck aus den Beobachtungen von \textsc{Bradley} reduziert hat.

Die jährliche Präzession der Sterne wurde nach den folgenden Formeln berechnet.
\begin{center}\small\renewcommand{\arraystretch}{1.18}\begin{tabular}{@{}l@{\hspace{8pt}}l@{}}
in Rektaszension $.\;.\;.\;.\;.\;.$ & $46''0395+20''0642\sin\alpha\operatorname{tang}\delta$ \\
in Deklination $.\;.\;.\;.\;.\;.\;.\;.\;.\;.$ & $20''0642\cos\alpha$ \\
\end{tabular}\end{center}
\noindent Diesen Zahlen liegen die jährliche Lunisolarpräzession in der Länge $50''388$, die jährliche Bewegung der Äquinoktialpunkte auf dem Äquator infolge der planetarischen Einwirkungen $0''1814$ und die mittlere Schiefe der Ekliptik für 1800, $23^\circ\;27'\;55''5$, zugrunde. Diese Bestimmungen stimmen sehr genau mit denen von \textsc{Bessel} überein. (Siehe das 19. Stück dieser Zeitschrift von 1814 [S. 561 dieses Bandes]).

Die äußere Gestaltung des Werks ist seines Inhalts durchaus würdig, doch ist der Druck etwas sparsamer ausgeführt als beim Verzeichnis von 1803. Die Rektaszensionen sind nur als Winkel genau angegeben, als Zeitangabe dagegen nur auf Minuten. Alle Sterne sind nach ihren Rektaszensionen geordnet, in 24 Stunden eingeteilt und innerhalb jeder einzelnen Stunde fortlaufend nummeriert. Bei allen Sternen mit eigenen Namen sind diese vor allem aus dem Verzeichnis von \textsc{Ulugh Beigh} beigefügt. Alle weiteren Bezeichnungen sind bei diesen Sternen dagegen weggeblieben. Das gefällt uns nicht, und wir glauben, dass alle Astronomen es unpraktisch finden werden, anstelle der ihnen vertrauten Buchstaben von \textsc{Bayer} und Zahlen von \textsc{Flamstead} nur die ihnen fast durchweg fremden arabischen Namen vorzufinden.

Am Schluss jeder einzelnen Stunde finden sich stets ausführliche Anmerkungen zu vielen der darin aufgeführten Sterne.

Nachdem der Katalog schon vollständig fertiggestellt und sogar gedruckt war, verglich ihn \textsc{Piazzi} noch mit den Verzeichnissen der Zodiakalsterne von \textsc{Zach} und \textsc{Barry}. Mit dem ersteren war die Übereinstimmung fast durchweg sehr groß. Bei den Mannheimer Deklinationen fanden sich dagegen häufigere und größere Unterschiede, von denen hier allerdings nur die größten angeführt sind, die ganze Bogenminuten erreichen.

